% From mitthesis package
% Version: 1.01, 2023/06/19
% Documentation: https://ctan.org/pkg/mitthesis
%
% The abstract environment creates all the required headers and footnote. 
% You only need to add the text of the abstract itself.
%
% Approximately 500 words or less; try not to use formulas or special characters
% If you don't want an initial indentation, do \noindent at the start of the abstract

Programmable integrated photonics (PIP) is emerging as a flexible and reconfigurable platform for realizing diverse optical functionalities on a single chip, enabling dynamic control of light propagation for communication, computing, and signal processing applications.
Among the available integration technologies, silicon photonics provides a scalable and cost efficient platform for high bandwidth, low power data transmission, making it a promising solution for next generation data center interconnects.
In particular, optical circuit switching enables direct optical connectivity for both scale-up and scale-out network architectures, avoiding optical–electrical–optical conversions.
However, the practical deployment of large scale programmable photonic switching fabrics is limited not only by hardware capabilities, but also by the scalable programming and control methodologies.

This thesis introduces a graph-based programming methodology that translates photonic circuits into weighted graphs, providing a mathematical abstraction of complex circuit fabrics.
Integrated circuit components and their connections are represented as graph elements with physical metrics, enabling high level programming, systematic optimization of hardware architectures and control methodologies.
Using this approach, routing and configuration algorithms are developed to support optical interconnects, multicasting, and circuit switching on a general purpose PIP processor, demonstrating that diverse networking functions can be realized on a single platform without manual intervention.
To address scalability challenges in large scale circuit fabrics, the graph-based methodology is extended to model and benchmark different topologies and routing strategies.
Incorporating physical layer metrics such as insertion loss, crosstalk, and signal-to-noise ratio allows quantitative evaluation of trade-offs and informs practical design guidance.
Finally, a software-defined control plane using model driven interfaces and gNMI is implemented to enable southbound communication, aligning Layer-1 optical resources with the SDN framework.

This work shows that a graph-based methodology provides hardware level programmability and design optimization for integrated photonic switches, while the developed software-defined control plane extends these capabilities to network level control and orchestration.
Together, these contributions establish a scalable foundation for integrating photonic circuit switching into AI-driven networking systems, supporting flexible, high performance, and manageable optical interconnects.
